\section{Discussion}
The results of this project are promising; the code finds an optimal orientation to print parts in given the constraints of compute time.  By utilizing methods that are meant for sparse data sets, the compute time is kept relatively low.  

However, this simulation only takes into account residual stress.  There are no considerations made for dimensional accuracy, print time, support strategy, or material usage.  As shown in the exhaust manifold, the direction of print gives greater z height.  This typically makes parts take longer to print.  Not only this, but printers are typically less accurate in the z direction (check that) so putting more print in that direction could have some tolerance stackup to contend with.  With this program, it is completely possible that excessive supports are created.  This not only kills print time, but increases material usage greatly.  If more supports are used, then more powder was used to print the supports, thus the print cost more.  In addition to this, there is no prohibition in the program for printing supports in inaccessible locations.  With LPBF, supports need to be removed via wire EDM, which requires an open path to cut.  If there is a set of supports inside of a tube for example, wire EDM isn't going to be able to cut them off.  In this case, more difficult methods such as machining or grinding could be used, but these are costly and less optimal.